\documentclass[11pt,a4paper]{article}
\usepackage{amsmath}
\usepackage{amssymb}
\usepackage{geometry}
\usepackage{graphicx}
\usepackage{hyperref}
\usepackage{listings}
\usepackage{xcolor}
\newcommand{\qed}{\hfill$\square$}

\geometry{margin=1in}

\title{Dugoff Tire Model for Mining Haul Trucks\\
\large Mathematical Derivation and Implementation}
\author{Mario Tilocca}
\date{\today}

\begin{document}

\maketitle

\begin{abstract}
This document presents the mathematical foundation of the Dugoff tire model for application to autonomous mining haul trucks operating in Western Australian desert conditions. The Dugoff model provides a computationally efficient representation of tire-ground interaction that captures friction saturation on loose surfaces while maintaining analytical tractability. We derive the complete model equations and provide parameter identification for the XCMG XDE320 electric dump truck (320-ton payload).
\end{abstract}

\section{Introduction}

The Dugoff tire model \cite{dugoff1970} is selected for mining haul truck simulation due to its key advantages:
\begin{itemize}
    \item Handles combined longitudinal and lateral slip (critical for cornering while accelerating/braking)
    \item Models friction saturation on loose surfaces (gravel, iron ore dust)
    \item Closed-form equations (no iterative solvers required)
    \item Captures smooth transition from linear to sliding regimes
\end{itemize}

\section{Fundamental Equations}

\subsection{Tire Force Model}

The Dugoff model expresses tire forces as:

\begin{equation}
F_x = C_x \sigma_x f(\lambda)
\end{equation}

\begin{equation}
F_y = C_y \sigma_y f(\lambda)
\end{equation}

where:
\begin{itemize}
    \item $F_x$ = longitudinal tire force [N]
    \item $F_y$ = lateral tire force [N]
    \item $C_x$ = longitudinal slip stiffness [N]
    \item $C_y$ = lateral slip stiffness [N]
    \item $\sigma_x$ = longitudinal slip ratio [-]
    \item $\sigma_y$ = lateral slip ratio [-]
    \item $f(\lambda)$ = friction saturation function
    \item $\lambda$ = friction utilization parameter
\end{itemize}

\subsection{Slip Ratio Definitions}

\textbf{Longitudinal Slip Ratio:}

\begin{equation}
\sigma_x = \frac{\omega R - V_x}{V_x}
\end{equation}

where $\omega$ is wheel angular velocity [rad/s], $R$ is tire effective rolling radius [m], and $V_x$ is vehicle longitudinal velocity [m/s].

\textbf{Sign Convention:}
\begin{itemize}
    \item $\sigma_x > 0$: Driving (wheel faster than vehicle)
    \item $\sigma_x < 0$: Braking (wheel slower than vehicle)
    \item $\sigma_x = 0$: Free rolling
\end{itemize}

\textbf{Lateral Slip Ratio:}

For small slip angles (typical in normal driving):

\begin{equation}
\sigma_y = \frac{V_y}{V_x} \approx \tan(\alpha) \approx \alpha
\end{equation}

where $V_y$ is lateral velocity at tire contact patch [m/s] and $\alpha$ is tire slip angle [rad].

\subsection{Friction Utilization Parameter}

The friction utilization parameter $\lambda$ determines whether the tire operates in the linear or friction-limited regime:

\begin{equation}
\lambda = \frac{\mu F_z}{2\sqrt{C_x^2 \sigma_x^2 + C_y^2 \sigma_y^2}}
\end{equation}

where:
\begin{itemize}
    \item $\mu$ = coefficient of friction (surface-dependent)
    \item $F_z$ = normal load on tire [N]
\end{itemize}

\textbf{Physical Interpretation:}
\begin{itemize}
    \item $\lambda \geq 1$: Linear regime (no sliding, forces proportional to slip)
    \item $\lambda < 1$: Nonlinear regime (friction-limited, approaching saturation)
    \item $\lambda \to 0$: Complete sliding (forces approach zero)
\end{itemize}

\subsection{Friction Saturation Function}

The saturation function $f(\lambda)$ is defined piecewise:

\begin{equation}
f(\lambda) = 
\begin{cases}
(2 - \lambda)\lambda = 2\lambda - \lambda^2, & \lambda < 1 \\
1, & \lambda \geq 1
\end{cases}
\end{equation}

This function ensures:
\begin{enumerate}
    \item Continuity at $\lambda = 1$: $\lim_{\lambda \to 1^-} f(\lambda) = 1$
    \item Smooth transition from linear to saturated regime
    \item Maximum value of $f(\lambda) = 1$ at $\lambda = 1$
    \item Monotonic decrease as $\lambda \to 0$
\end{enumerate}

\section{Force Saturation Proof}

\textbf{Theorem:} The Dugoff model guarantees that total tire force never exceeds available friction:

\begin{equation}
\sqrt{F_x^2 + F_y^2} \leq \mu F_z
\end{equation}

\textbf{Proof:}

Start with the total force magnitude:
\begin{align}
\sqrt{F_x^2 + F_y^2} &= \sqrt{(C_x \sigma_x f(\lambda))^2 + (C_y \sigma_y f(\lambda))^2} \\
&= f(\lambda) \sqrt{C_x^2 \sigma_x^2 + C_y^2 \sigma_y^2}
\end{align}

For $\lambda < 1$ (friction-limited regime), substitute $f(\lambda) = 2\lambda - \lambda^2$ and $\lambda$ from Eq. (5):

\begin{align}
\sqrt{F_x^2 + F_y^2} &= (2\lambda - \lambda^2) \sqrt{C_x^2 \sigma_x^2 + C_y^2 \sigma_y^2} \\
&= (2\lambda - \lambda^2) \cdot \frac{\mu F_z}{2\lambda} \\
&= \left(2 - \lambda\right) \frac{\mu F_z}{2} \\
&< \mu F_z \quad \text{(since } 0 < \lambda < 1\text{)}
\end{align}

For $\lambda \geq 1$ (linear regime), $f(\lambda) = 1$, so the tire is not friction-limited and the inequality holds by construction. \qed

\section{Parameter Identification}

\subsection{Tire Slip Stiffness}

For 40.00R57 mining tires on XCMG XDE320:

\textbf{Base Values at Reference Load:}
\begin{align}
C_{x,\text{base}} &= 280{,}000 \text{ N} \\
C_{y,\text{base}} &= 220{,}000 \text{ N} \\
F_{z,\text{ref}} &= 800{,}000 \text{ N (800 kN)}
\end{align}

\textbf{Load-Dependent Scaling:}

Slip stiffness increases with normal load following an empirical square-root relationship:

\begin{align}
C_x(F_z) &= C_{x,\text{base}} \left(\frac{F_z}{F_{z,\text{ref}}}\right)^{0.5} \\
C_y(F_z) &= C_{y,\text{base}} \left(\frac{F_z}{F_{z,\text{ref}}}\right)^{0.5}
\end{align}

This nonlinear scaling captures the stiffening effect of increased contact patch pressure.

\subsection{Friction Coefficients}

Surface-dependent friction coefficients for Pilbara mining conditions:

\begin{table}[h]
\centering
\begin{tabular}{lcc}
\hline
\textbf{Surface Type} & $\mu_{\text{peak}}$ & $\mu_{\text{slide}}$ \\
\hline
Dry pavement (smooth haul road) & 0.85 & 0.75 \\
Gravel compact (baseline) & 0.72 & 0.65 \\
Gravel loose (after grading) & 0.55 & 0.48 \\
Iron ore dust (dry) & 0.45 & 0.38 \\
Iron ore dust (wet) & 0.30 & 0.25 \\
Mud (wet season) & 0.25 & 0.20 \\
Sand (loose) & 0.40 & 0.35 \\
\hline
\end{tabular}
\caption{Friction coefficients by surface type}
\label{tab:friction}
\end{table}

\textbf{Note:} Peak friction $\mu_{\text{peak}}$ applies at low slip ratios ($|\sigma_x| < 0.05$). Slide friction $\mu_{\text{slide}}$ applies when wheels are locked ($|\sigma_x| \approx 1$). The Dugoff model naturally interpolates between these values.

\subsection{Normal Load Distribution}

For XCMG XDE320 (empty weight = 218 tons, GVW = 538 tons):

\textbf{Static Load Distribution:}

\begin{align}
F_{z,\text{front}} &= \frac{mg \cdot d_{\text{rear}}}{L} \\
F_{z,\text{rear}} &= \frac{mg \cdot d_{\text{front}}}{L}
\end{align}

where:
\begin{itemize}
    \item $m$ = vehicle mass [kg]
    \item $g = 9.81$ m/s$^2$
    \item $L = 6.3$ m (wheelbase)
    \item $d_{\text{front}}$ = distance from CoG to front axle [m]
    \item $d_{\text{rear}}$ = distance from CoG to rear axle [m]
\end{itemize}

\textbf{Dynamic Load Transfer (Longitudinal):}

During acceleration or braking, load transfers between axles:

\begin{align}
\Delta F_z &= \frac{m a_x h_{\text{CoG}}}{L} \\
F_{z,\text{front}} &= F_{z,\text{front,static}} - \Delta F_z \\
F_{z,\text{rear}} &= F_{z,\text{rear,static}} + \Delta F_z
\end{align}

where $a_x$ is longitudinal acceleration [m/s$^2$] and $h_{\text{CoG}}$ is center of gravity height [m].

\textbf{Sign Convention:}
\begin{itemize}
    \item $a_x > 0$: Acceleration (weight transfers to rear)
    \item $a_x < 0$: Braking (weight transfers to front)
\end{itemize}

\newpage
\section{Numerical Implementation}

\subsection{Algorithm}

For each tire at each timestep:

\begin{enumerate}
    \item Compute slip ratios $\sigma_x$, $\sigma_y$ from kinematics
    \item Compute normal load $F_z$ from load distribution
    \item Compute slip stiffness $C_x(F_z)$, $C_y(F_z)$ from load scaling
    \item Compute friction utilization $\lambda$ from Eq. (5)
    \item Compute saturation function $f(\lambda)$ from Eq. (6)
    \item Compute tire forces $F_x$, $F_y$ from Eqs. (1)-(2)
\end{enumerate}

\subsection{Numerical Stability}

\textbf{Division by Zero Protection:}

When computing slip ratios at low speeds:

\begin{equation}
\sigma_x = \frac{\omega R - V_x}{\max(V_x, V_{\text{min}})}
\end{equation}

where $V_{\text{min}} = 0.5$ m/s prevents division by zero.

\textbf{Epsilon in Lambda Calculation:}

To avoid division by zero when slip is zero:

\begin{equation}
\lambda = \frac{\mu F_z}{2\sqrt{C_x^2 \sigma_x^2 + C_y^2 \sigma_y^2} + \epsilon}
\end{equation}

where $\epsilon = 10^{-6}$ N.

\section{Expected Behavior}

\subsection{Traction-Limited Acceleration}

For XCMG XDE320 loaded (538 tons) on gravel compact ($\mu = 0.72$):

\begin{align}
F_{z,\text{rear}} &\approx 0.75 \times 538{,}000 \times 9.81 = 3{,}956{,}000 \text{ N} \\
F_{x,\text{max}} &= \mu F_{z,\text{rear}} = 0.72 \times 3{,}956{,}000 = 2{,}848{,}000 \text{ N} \\
F_{x,\text{motor}} &= \frac{T_{\text{motor}}}{R} = \frac{145{,}000}{1.93} = 75{,}100 \text{ N}
\end{align}

Result: Motor force $\ll$ traction limit, so truck is \textbf{power-limited}, not traction-limited.

\subsection{Emergency Braking}

For same truck on wet iron ore dust ($\mu = 0.30$):

\begin{align}
F_{x,\text{max}} &= 0.30 \times 538{,}000 \times 9.81 = 1{,}585{,}000 \text{ N} \\
F_{x,\text{brake}} &= \frac{T_{\text{brake}}}{R} = \frac{180{,}000}{1.93} = 93{,}300 \text{ N}
\end{align}

Result: Still power-limited, but much closer to traction limit. On loose gravel ($\mu = 0.55$), wheels would lock.

\section{Validation}

\subsection{Unit Tests}
\begin{itemize}
    \item Verify $f(\lambda)$ continuity at $\lambda = 1$
    \item Confirm $\sqrt{F_x^2 + F_y^2} \leq \mu F_z$ for random inputs
    \item Test load transfer during acceleration/braking
\end{itemize}

\subsection{Integration Tests}
\begin{itemize}
    \item Straight-line acceleration on various surfaces
    \item Emergency braking with ABS simulation
    \item Cornering at limit friction
    \item Combined braking and cornering (friction circle)
\end{itemize}

\section{Conclusion}

The Dugoff tire model provides a computationally efficient framework for simulating mining haul truck dynamics on loose surfaces. Its closed-form equations and guaranteed friction saturation make it suitable for real-time simulation and control algorithm development.

Future work includes:
\begin{itemize}
    \item Validation against field data from XCMG XDE320
    \item Extension to 3D terrain (grades, banking)
    \item Thermal effects on friction coefficients
    \item Comparison with Pacejka Magic Formula
\end{itemize}

\begin{thebibliography}{9}
\bibitem{dugoff1970}
Dugoff, H., Fancher, P. S., and Segel, L.,
\textit{An Analysis of Tire Traction Properties and Their Influence on Vehicle Dynamic Performance},
SAE Technical Paper 700377, 1970.

\bibitem{pacejka2012}
Pacejka, H.,
\textit{Tire and Vehicle Dynamics}, 3rd Edition,
Butterworth-Heinemann, 2012.

\bibitem{rajamani2012}
Rajamani, R.,
\textit{Vehicle Dynamics and Control}, 2nd Edition,
Springer, 2012.
\end{thebibliography}

\end{document}